\section{模式判断需要怎样的联合预言}\label{sec:3-8}

在低能核结构高自旋谱学研究中，研究者常常依据个别能级比值接近某种理想转子或振子规律，或依据观测到一条相对增强的带间跃迁，便赋予实验能带特定的模式标签（如“手征双重带”、“横向摆动带”或“八极反射不对称带”）。然而，正如第\ref{chap:2}章第\ref{sec:2-8}节的综合比较所指出的，不同物理模型由于包含若干可调参数，在低激发能区往往表现出显著的拟合简并。若仅凭能谱局域拟合便断言某种运动模式已确立，容易将参数自由度带来的数值符合误认为物理机制的确证。

为了克服这一局限，本节旨在探讨如何将模式候选转化为具有明确前向预言、最强替代解释和判决性检验路径的研究问题：首先梳理理论模型的前向预言架构与参数输入；进而分析竞争解释在低能区的参数重叠与不可通约性，明确为何不能设立通用的唯象绝对阈值；接着建立包含误差协方差的数据分层架构与留出检验准则；最后基于已有文献基线，提出三个具体可证伪的实证研究候选方案（$^{98}\mathrm{Mo}$ 弱 $E3$ 跃迁、$^{102}\mathrm{Ru}$ 低能偶极聚集与 $A \approx 130$ 低自旋取向进动），为后续全书跨质量区实测分析与实验统计反演提供方法论参照。\cite{lawrie_2020_Tiltedprecession,frauendorf_2014_Transversewobbling,nomura_2022_Questioningwobbling,lv_2022_Evidencewobbling,karmakar_2025_98mo_octupole,saracino_2026_102ru_dipole}

\subsection{模式的前向预言结构与模型空间划分}\label{sec:3-8-1}

一个自洽的核结构物理模型，不是事后对实验能级的唯象拟合公式，而是基于明确物理假定对多维物理量的前向映射系统。无论采用全微观自洽理论还是具有明确假定的唯象代数模型，建立一个模式的前向理论预言，均须明确以下四个核心要素：
\begin{enumerate}
\item \textbf{选定的动力学自由度划分}：明确模型保留的主导动力学变量（如转动欧拉角取向坐标、玻尔四极形状坐标 $\beta, \gamma$ 或准粒子产生湮灭算符），以及被冻结或消去的环境空间（如核芯高能虚激发或高阶多极形变）。
\item \textbf{守恒量子数与对称性约束}：明确态波函数具有的确定总角动量与宇称 $I^\pi$，以及所保持的特定离散对称性类别（如绕主轴旋转 $\pi$ 的签名量子数；时间反演作为反幺正操作，其诱导的代数约束不同于普通的幺正宇称本征标签）。
\item \textbf{势能面与动力学输入}：明确底层势能曲面（微观自洽势面或唯象变形势阱）、转动惯量张量 $\mathcal{J}_k(\omega)$、微观配对能隙 $\Delta$ 以及经模型空间截断重整化的有效跃迁算符。
\item \textbf{拟合参数与前向预言的严格划分}：必须清晰标明哪些实验量被用于固定模型的自由参数（如基态带转动常数或单粒子能隙），模型的前向预言能力严格体现在那些未参与拟合的独立物理量上。
\end{enumerate}

需要指出，模型输出的内禀波函数成分分布或半经典相空间角动量进动轨迹（如倾斜进动平面或相空间分离线），仅是特定模型表象下的中间数学表达，其本身并不是直接可观测的实验物理量。将模型标签与实验观测画等号，在认识论上混淆了理论表象与客观物理测量。

\subsection{竞争解释的参数重叠与阈值不可通约性}\label{sec:3-8-2}

在低自旋区间，不同的微观激发机制常常展现出高度重叠的实验信号特征：
\begin{itemize}
\item 组态相近的单粒子准粒子对顺排带、具有相同宇称的签名伴带、大振幅 $\gamma$ 软核表面激发、粒子-转子倾斜进动以及转子摆动运动，在低激发区均可能给出近乎相同的转动能级间隔序列，甚至同样表现出带间电四极跃迁。
\item 在传统唯象分析中，常有人试图设定通用的判定阈值（例如规定多极混合比量级必须满足 $|\delta| > 1$、或带间 $E2$ 分支大于 $50\%$、或二声子能级必须恰好等于单声子能级的两倍）。这种通用阈值假定缺乏理论自洽性：不同核素的单粒子组态、核芯形变刚度以及科里奥利反演点各不相同，不存在放之四海而皆准的绝对分界线。
\item 进一步，角分布两点各向异性比值 $R_{\mathrm{ac}}$ 在勒让德多项式展开下导出关于多极混合比 $\delta$ 的二次代数方程。根据判别式的符号与实验误差边界，数学上可能出现两个实根分支、单个极小解或在特定置信区间内无物理实解。在缺乏独立线偏振度测量进行有效限制时，必须在分析中客观保留相容的 $\delta$ 分支，不能依据特定理论偏好人为删除任一分支；同时，保留多分支拟合解亦不等于宣告这些解在物理机制上先验等价。
\item 此外，光子多极混合比 $\delta(E2/M1)$ 严格由第\ref{sec:5-9}节给出的电四极与磁偶极多极辐射振幅之比定义，包含了跃迁能量因子、辐射场波长与标准电磁相位。当磁偶极跃迁矩阵元因几何角动量取向或单粒子轨道禁戒而导致分母较小时，即使绝对 $B(E2)$ 强度并未发生显著的集体增强，数值上亦可呈现出较大的 $|\delta|$ 值。因此，绝对折合跃迁几率 $B(E2)$ 与相对分支比或混合比具有本质不同的物理信息，二者不可相互替代。
\end{itemize}

\subsection{独立观测量的协方差与四层数据架构}\label{sec:3-8-3}

在评估实验证据对理论模型的支持度时，必须审慎考察观测量之间的统计独立性：
\begin{enumerate}
\item \textbf{寿命与跃迁强度的统计协方差}：对于由同一激发母态退激产生的多条伽马跃迁分支，实验提取的折合跃迁几率 $B(\sigma\lambda)$ 反比于母态平均寿命 $\tau$，且正比于该分支的部分跃迁几率：
\begin{equation}\label{eq:3-8-b-rate-relation}
B(\sigma\lambda) \propto \lambda_{\gamma} = \frac{b_\gamma}{\tau (1 + \alpha_{\mathrm{IC}})} ,
\end{equation}
其中 $b_\gamma$ 为光子分支比，$\alpha_{\mathrm{IC}}$ 为总内转换系数。式\eqref{eq:3-8-b-rate-relation}表明，寿命越短则对应的绝对跃迁几率越大。由于同一能级退激的多条分支共用同一个实测母态寿命 $\tau$ 以及总计数探测效率标定，由该母态派生出的多个跃迁强度 $B(\sigma\lambda)$ 之间存在不可忽略的统计协方差（Covariance）。不能将共用同一套数据采集的多个衍生强度当作相互独立的统计互证。
\item \textbf{能谱调参不属于留出验证}：利用同一核素的低自旋能级拟合模型参数后，再次用该模型计算同一能谱的吻合程度，仅属于自洽闭环拟合，不具备独立的检验效力。自洽的模型验证必须将已调参数严格固定，并在完全未参与参数调整的“留出观测空间”中接受检验。
\item \textbf{四层数据架构}：为了避免证据层级的混淆，研究中必须严格区分四类数据载荷：
\begin{itemize}
\item \textbf{原始测量量 (Raw Data)}：探测器原始脉冲计数、符合矩阵、原始各向异性角分布计数；
\item \textbf{采用物理量 (Adopted Properties)}：经探测器响应反演后的能级能量、激发态半衰期 $T_{1/2}$；
\item \textbf{派生跃迁量 (Derived Observables)}：结合分支比与转换系数计算出的绝对折合跃迁几率 $B(E2), B(M1)$ 及拟合提取的混合比 $\delta$；
\item \textbf{理论模型输出 (Model Quantities)}：模型本征能量、组态混合矩阵元与相空间角动量轨迹。
\end{itemize}
\end{enumerate}

在进行模型与实验数据的加权比对时，残差向量 $\mathbf{r} = \mathbf{y}_{\mathrm{exp}} - \mathbf{y}_{\mathrm{model}}$ 经实验协方差矩阵逆矩阵的平方根加权后（$\mathbf{V}^{-1/2} \mathbf{r}$），方构成统计正交残差。在频率学统计框架下，基于已知实验中心值与公布协方差矩阵构造的加权距离 $\chi^2 = \mathbf{r}^T \mathbf{V}^{-1} \mathbf{r}$ 具有自洽性，并不依赖贝叶斯先验假设；但严禁跨越非嵌套模型滥用拟合优度判定理论机制的绝对优劣，更为细致的多维非线性误差传播与统计假设检验将在第\ref{sec:5-10}节展开。

\subsection{三个可证伪的实证研究候选}\label{sec:3-8-4}

基于上述方法论标准，本小节提出三个具有确凿文献原件基线、物理机制明确且具备可操作判决路径的实证研究候选方案。

\paragraph{候选一：$^{98}\mathrm{Mo}$ 弱 $E3$ 跃迁与 $B(E3)/B(E2)$ 增强机制的实验判决}
在 $A \approx 100$ 质量区的富中子偶偶核 $^{98}\mathrm{Mo}$（$Z=42, N=56$）中，Karmakar 等人在近期高自旋研究中观测到了连接负宇称带与正宇称带的两条弱跃迁候选。\cite{karmakar_2025_98mo_octupole}
\begin{itemize}
\item \textbf{已有测量与能级结构}：实验在能级纲图中建立了 $1833.3\,\mathrm{keV}$（$5^- \to 2_1^+$）与 $1585.8\,\mathrm{keV}$（$7^- \to 4_1^+$）两条弱跨带跃迁；由于计数统计受限，实验中未能提取出这两条弱跃迁可靠的角关联 DCO 比值与偏振度，由始末态自旋宇称差确定电八极 $E3$ 为其最低允许的多极阶数。
\item \textbf{物理缺口与最强竞争解释}：目前负宇称母态（$5^-$ 与 $7^-$）的绝对寿命尚未测定。显微分析表明，测量到的较小分支比强度存在两种性质不同的竞争解释：其一为分子八极矩阵元 $\langle || \mathcal{M}(E3) || \rangle$ 表现出集体增强；其二为分母上的带内参考 $E2$ 跃迁因初末态形状变化或多极阻碍而受到压低。原文献在特定重离子反应道下未提取出可用线形，并不代表该寿命原则上不可测。
\item \textbf{判决性实验检验路径}：需要依据探测器响应反演与复杂的级联侧馈送（side feeding）模型，进行寿命灵敏度预研究；针对纳秒至皮秒的不同时间窗口，评估反冲距离多普勒位移法（RDM）或连续多普勒频移衰减法（DSAM）的可行性，测定该负宇称态的绝对寿命 $\tau$。通过绝对寿命将相对分支比解耦为绝对折合跃迁几率 $B(E3)$ 与 $B(E2)$，是区分“八极集体增强”与“四极形状受阻”的关键观测量。
\end{itemize}

\paragraph{候选二：$^{102}\mathrm{Ru}$ 低能偶极聚集与八极共振结构检验}
在近球形核素 $^{102}\mathrm{Ru}$（$Z=44, N=58$）中，Saracino 等人利用极化光子束核共振荧光（NRF）技术，在低激发能区观测到偶极跃迁强度的聚集分布。\cite{saracino_2026_102ru_dipole}
\begin{itemize}
\item \textbf{已有测量与选择定则界限}：NRF 实验测量给出的直接观测量是基态跃迁能级衰变宽度 $\Gamma_0$（或散射截面积分 $I_s \propto \Gamma_0^2/\Gamma$）。实验在 $2.95\,\mathrm{MeV}$ 处测得明确的电偶极跃迁，折合跃迁几率达到 $B(E1)\!\uparrow = 0.016\,e^2\mathrm{fm}^2$，并在 $4.5\sim 5.5\,\mathrm{MeV}$ 能区观测到总强度约为 $0.1\,e^2\mathrm{fm}^2$ 的强偶极聚集。需要特别澄清：由于基态为 $0^+$，测得的基态跃迁 $0^+ \to 1^-$ 严格遵循电偶极 $E1$ 跃迁选择定则，单光子吸收过程中角动量守恒限制了 $0^+$ 态不能直接通过电八极 $E3$ 算符激发到同一 $1^-$ 态（角动量三角定则要求 $|0 - 1| \le \lambda \le 0 + 1 \implies \lambda = 1$）。因此，不能将对 $1^-$ 态的 $E1$ 测量等同于对该态直接进行 $E3$ 测量。
\item \textbf{物理缺口与竞争模型}：争论的焦点在于，低能 $1^-$ 态的聚集究竟代表二声子八极-四极相干耦合构型 $[2^+ \otimes 3^-]_{1^-}$ 的投射，还是由非相干准粒子微观极化产生的局域共振。
\item \textbf{判决性实验检验路径}：检验八极声子模型假定，需要利用在束谱学测量该 $1^-$ 态向八极声子带成员 $3_1^-$ 态的 $E2$ 连接跃迁（$1^- \to 3_1^-$），或向基态带 $2_1^+$ 态的跨带跃迁，并结合基态到 $3_1^-$ 态独立的 $0^+ \to 3_1^-$ 库仑激发 $B(E3)$ 测量，检验微观势面预言的八极多声子耦合矩阵元。
\end{itemize}

\paragraph{候选三：$A \approx 130$ 奇 $A$ 核低自旋取向进动模式的判决路径}
在典型的过渡区奇 $A$ 核（以 $^{135}\mathrm{Pr}$ 与 $^{133}\mathrm{La}$ 为代表）中，围绕低自旋负宇称能带的动力学机制存在不同的理论模型描述。\cite{lawrie_2020_Tiltedprecession,frauendorf_2014_Transversewobbling,nomura_2022_Questioningwobbling,lv_2022_Evidencewobbling}
\begin{itemize}
\item \textbf{测量结果与分析差异}：不同实验组基于各自采集的数据与分析程序，得出了不同的多极混合比解。Sensharma 等人在角分布分析中选取了较大混合比分支，并结合冻结准粒子取向的粒子-转子模型（PRM）解释为横向摆动；而 Lv 等人结合康普顿散射线偏振度实测数据，指出同一批数据同样容许以磁偶极为主导的小根解（对 $747.3\,\mathrm{keV}$ 跃迁提取出 $\delta = -0.47^{+0.09}_{-0.22}$），并在准粒子转子（QTR）模型中将其描述为准粒子角动量向中间轴顺排的倾斜进动。这两类模型与 Nomura 等人提出的 IBFM 软核表面激发模型并不对应于全域互斥的排他标签，而是反映了不同自由度截断下的理论视角。
\item \textbf{判决性实验检验路径}：不能人为忽视角分布分析中可能存在的多解分支，而应设计更为严谨的实验判决：(a) 利用离散高纯锗探测阵列的多角度符合测量，构建包含探测器立体角卷积的联合似然函数，检验能否在统计置信度内对 $\delta$ 的物理分支形成有效约束；(b) 进行探测器响应与寿命灵敏度预研究，评估反冲距离测量方法的可行性，测定关键带间跃迁的绝对折合跃迁几率 $B(E2)$，将其与各模型给出的理论预测区间进行后验比对；(c) 结合高自旋区准粒子退耦顺排测量，检验角动量进动几何随自旋演化的微观响应。
\end{itemize}

\subsection{跨质量区迁移准则与全书判据总结}\label{sec:3-8-5}

图\ref{fig:3-8-joint-workflow}概括了从理论模型空间出发，经由多维互补信息通道，最终落地于实证研究候选的物理闭环。

\begin{figure}[htbp]
\centering
\includegraphics[width=0.88\textwidth]{Img/joint-predictions-workflow.pdf}
\caption{模式判断需要怎样的联合预言：理论前向预言架构、多维互补信息通道与可证伪实证研究候选闭环示意图（由独占脚本 \texttt{system/scripts/a130\_joint\_predictions\_figure.py} 生成）。图示强调理论模型假定划分、互补观测量留出检验（含分支比协方差）、以及三个具体实证候选路径的逻辑对应。}\label{fig:3-8-joint-workflow}
\end{figure}

当将某一特定核区（如 $A \approx 130$ 区）总结出的模式机制推广迁移至其它质量区时，必须严格遵守微观自洽的物理迁移准则：
\begin{itemize}
\item 必须逐一核查目标核区的微观费米面轨道占据情况、自旋-轨道伴对轨道的相对劈裂、质子-中子相互作用的有效强度、自洽平均场势能面的三轴刚度、以及反应布居通道的自旋输入特征。
\item 严禁仅因某个质量数区间与已知变形核相近，或因本地文献库尚未检索到相关报道，便轻率断言其“首次确证”或“未曾发现”。原子核的多体动力学状态严格由其微观壳结构与残余相互作用决定，不随人为的文献分类而转移。
\end{itemize}

综上所述，将一个理论模式确立为真实客观的核物理激发，绝非依赖孤立的能谱拟合或人为的模式命名，而是依赖于一组具有互补信息维度、相互自洽的联合预言链条。通过建立包含绝对寿命、多角度偏振分析、高自旋顺排演化以及严格考量统计协方差的多维观测体系，并面向具体实证候选开展判决性测量，才能真正打破模型解释简并，将低能核结构高自旋谱学研究推进至自洽严谨的科学高度。
